% from ega2-1-fr.tex
\oldpage[II]{201}
on peut prendre $X'=X$), alors l'homomorphisme $\alpha$
(\hyperref[II.8.14.5.4-fr]{8.14.5.4}) est \textbf{(TN)}-\emph{bijectif} pour
tout $\mathcal{S}$-Module gradué quasi-cohérent $\mathcal{M}$ vérifiant la
condition \textbf{(TF)}.
\end{remark}

\begin{remark}[8.14.12]
\label{II.8.14.12-fr}
La situation décrite dans (\hyperref[II.8.14.4-fr]{8.14.4}) est un cas
particulier de la suivante. Soient $X$ un espace annelé, $\mathcal{S}$ une
$\mathcal{O}_X$-Algèbre graduée (à degrés positifs et négatifs)~; supposons
qu'il existe un entier $d>0$ tel que $\mathcal{S}_d$ et
$\mathcal{S}_{-d}$ soient \emph{inversibles}, l'homomorphisme canonique
\begin{equation}
\label{II.8.14.12.1-fr}
  \mathcal{S}_d\otimes_{\mathcal{O}_X}\mathcal{S}_{-d}
  \to\mathcal{O}_X
\tag{8.14.12.1}
\end{equation}
étant un \emph{isomorphisme} (de sorte que $\mathcal{S}_{-d}$ s'identifie à
$\mathcal{S}_d^{-1}$). On dit alors que la $\mathcal{O}_X$-Algèbre graduée
$\mathcal{S}$ est \emph{périodique de période $d$}. Cette dénomination
provient de la propriété suivante~: \emph{sous les hypothèses précédentes,
pour tout $\mathcal{S}$-Module gradué $\mathcal{F}$, l'homomorphisme
canonique}
\begin{equation}
\label{II.8.14.12.2-fr}
  \mathcal{S}_d\otimes\mathcal{F}_n\to\mathcal{F}_{n+d}
\tag{8.14.12.2}
\end{equation}
\emph{est un isomorphisme pour tout $n\in\mathbf{Z}$.} En effet, la question
est locale sur $X$ et on peut donc supposer que $\mathcal{S}_d$ possède une
section \emph{inversible} $s$ au-dessus de $X$, son inverse $s'$ étant une
section de $\mathcal{S}_{-d}$. L'homomorphisme
$\mathcal{F}_{n+d}\to\mathcal{S}_d\otimes\mathcal{F}_n$ qui, à toute section
$z\in\Gamma(U,\mathcal{F}_{n+d})$ fait correspondre la section
$(s|U)\otimes(s'|U)z$ de $\mathcal{S}_d\otimes\mathcal{F}_n$ sur $U$, est
alors réciproque de (\hyperref[II.8.14.12.2-fr]{8.14.12.2}), d'où notre
assertion. On en déduit pour tout $k\in\mathbf{Z}$ un isomorphisme canonique
\[
  (\mathcal{S}_d)^{\otimes k}\otimes\mathcal{F}_n
  \xrightarrow{\sim}\mathcal{F}_{n+kd}.
\]
Par suite, \emph{un $\mathcal{S}$-Module gradué $\mathcal{F}$ est connu
lorsqu'on connaît les $\mathcal{S}_0$-Modules $\mathcal{F}_i$
$(0\leqslant i\leqslant d-1)$ et les homomorphismes canoniques}
\[
  \mathcal{S}_i\otimes\mathcal{F}_j\to\mathcal{F}_{i+j}
  \qquad\text{pour }0\leqslant i,j\leqslant d-1
\]
(en posant
$\mathcal{F}_{i+j}=\mathcal{S}_d\otimes_{\mathcal{S}_0}
\mathcal{F}_{i+j-d}$ lorsque $i+j\geqslant d$). Bien entendu, pour que ces
homomorphismes définissent bien une structure de $\mathcal{S}$-Module sur la
somme directe des $(\mathcal{S}_d)^{\otimes k}\otimes\mathcal{F}_i$
$(k\in\mathbf{Z},\ 0\leqslant i\leqslant d-1)$, ils doivent satisfaire à des
conditions d'associativité que nous n'expliciterons pas.

Dans le cas où $d=1$ (qui est celui considéré dans
(\hyperref[subsection:II.3.3-fr]{3.3})) on peut donc dire que la catégorie des
$\mathcal{S}$-Modules gradués (resp. quasi-cohérents si $X$ est un préschéma
et $\mathcal{S}$ est quasi-cohérente) est \emph{équivalente} à la catégorie
des $\mathcal{S}_0$-Modules quelconques (resp. quasi-cohérents)~; c'est dans
ce sens qu'on peut considérer les développements de ce numéro comme
généralisant ceux du \S~\hyperref[section:II.3-fr]{3}. En outre, on voit que,
sous des conditions de finitude convenables, les résultats de ce numéro
(joints à (\hyperref[II.8.14.11-fr]{8.14.11})) ramènent dans une certaine
mesure l'étude des Algèbres graduées quasi-cohérentes sur un préschéma, et des
Modules gradués «~mod. \textbf{(TN)}~» sur de telles Algèbres, à l'étude du
cas particulier où les Algèbres considérées sont \emph{périodiques} (et où
la condition \textbf{(TN)} pour $\mathcal{M}$
(\hyperref[II.3.4.2-fr]{3.4.2}) implique donc $\mathcal{M}=0$).
\end{remark}

\begin{remark}[8.14.13]
\label{II.8.14.13-fr}
Sous les hypothèses de (\hyperref[II.8.14.1-fr]{8.14.1}), soit $d$ un entier
$>0$~; on a défini un $Y$-isomorphisme canonique $h$ de $X$ sur
$X^{(d)}=\operatorname{Proj}(\mathcal{S}^{(d)})$
(\hyperref[II.3.1.8-fr]{3.1.8}). Pour tout

